\section{Outer code instantiation I: a dense log-memory convolutional outer}
\label{sec:outer-dense-logmem}

This section instantiates the outer code \(\Cone\) by a simple, fully explicit ensemble:
a terminated, feedforward convolutional encoder with memory \(M=\Theta(\log n)\) and dense taps
within the memory window. We prove the only two outer properties used by the serial concatenation
framework:
\begin{enumerate}
\item \textbf{(O1) Logarithmic distance:} \(d_{\min}(\Cone)\ge \alpha\log n\) with high probability over the encoder design.
\item \textbf{(O2) Spectral growth:} \(\enum_{\Cone,w}\le \poly(n)\,\beta^w\) for a constant \(\beta>1\) (independent of \(n\))
      on the weight range used later in the first-moment bound.
\end{enumerate}
The construction is encoding-friendly in the sense that it is a single banded (width \(M+1\)) generator
with no feedback; its work is \(O(nM)=O(n\log n)\).

\subsection{Construction (banded diagonal feedforward convolution)}
\label{sec:outer-dense-construction}

Fix an outer rate \(R_o\in(0,1)\). Let \(k:=\lfloor R_o n\rfloor\) and let \(r:=1/R_o\) be the number of output streams
(so \(n=rk\), up to rounding).
Let \(M:=\lceil c\log_2 n\rceil\) for a constant \(c\ge 1\) chosen later.

Given information bits \(x_1,\dots,x_k\in\F_2\), we use standard zero-termination:
append \(M\) zeros \(x_{k+1}=\cdots=x_{k+M}=0\) and define \(x_t:=0\) for all \(t\le 0\).

For each time \(t\in[k]\) and stream \(j\in[r]\), sample an independent uniform tap vector
\(g_t^{(j)}\leftarrow \F_2^{M+1}\), and output
\begin{equation}
\label{eq:outer-dense-encoder}
y_t^{(j)} \;:=\; \big\langle g_t^{(j)},\, (x_t,x_{t-1},\dots,x_{t-M})\big\rangle \in \F_2.
\end{equation}
Flatten \(\bigl(y_t^{(1)},\dots,y_t^{(r)}\bigr)_{t=1}^k\) into a length-\(n\) outer codeword \(y\in\F_2^n\).
This defines a linear outer code \(\Cone\subseteq\F_2^n\) with generator matrix \(\Gone\in\F_2^{k\times n}\)
that is banded diagonal with bandwidth \(M+1\) (dense within the band, no feedback).

\paragraph{Outer spectrum notation.}
For \(w\in\{0,1,\dots,n\}\), let
\[
\enum_{\Cone,w} \;:=\; \bigl|\{c\in\Cone:\ \wt(c)=w\}\bigr|
\]
denote the (outer) weight spectrum.

\subsection{Spectrum bound (O2) via a trellis transfer matrix}
\label{sec:outer-dense-spectrum}

The encoder admits a standard trellis with \(2^M\) states (the last \(M\) input bits),
and constant out-degree \(2\) (the next input bit is \(0\) or \(1\)).
Each trellis edge produces \(r\) output bits, hence has Hamming weight in \(\{0,1,\dots,r\}\).

\begin{lemma}[Outer spectrum bound (trellis form)]
\label{lem:outer-dense-spectrum}
There exists a constant \(\beta>1\) (depending only on \(r\) and \(c/\alpha\) as folded below) such that
\[
\enum_{\Cone,w} \;\le\; \poly(n)\,\beta^w
\]
for all \(w\ge \alpha\log n\). (This is the range used by the framework, since \(d_{\min}(\Cone)\ge \alpha\log n\).)
\end{lemma}

\begin{proof}
Let \(T(z)\in\R[z]^{2^M\times 2^M}\) be the weighted adjacency matrix of the trellis, defined by
\[
T(z)_{s,s'} \;:=\; \sum_{e:s\to s'} z^{\wt(e)},
\]
where the sum is over edges \(e\) from state \(s\) to \(s'\) and \(\wt(e)\in\{0,\dots,r\}\) is the edge output weight.
Each state has exactly two outgoing edges, and each edge contributes at most \(z^r\), so every row-sum is bounded by
\(1+z^r\). Hence for all \(z\ge 0\),
\[
\rho\!\bigl(T(z)\bigr)\le 1+z^r
\]
by the maximum row-sum bound on the spectral radius.

Let \(W(z):=\sum_{w\ge 0}\enum_{\Cone,w} z^w\) be the (ordinary) weight enumerator.
Standard transfer-matrix reasoning gives
\[
W(z)\;\le\; \mathbf{1}^\top T(z)^k \mathbf{1} \;\le\; 2^M\cdot (1+z^r)^k
\]
for all \(z\ge 0\), where \(2^M\) accounts for the number of possible start/end states (termination only helps).

Extract coefficients using \(\enum_{\Cone,w}\le W(z)/z^w\) for any \(z>0\):
\[
\enum_{\Cone,w}\le 2^M\cdot \frac{(1+z^r)^k}{z^w}.
\]
Choose \(z\) so that \(z^r = w/n\) (i.e., \(z=(w/n)^{1/r}\)). Then \((1+z^r)^k=(1+w/n)^k\le (1+w/n)^n\le e^w\),
and thus
\[
\enum_{\Cone,w}\le 2^M\cdot e^w \cdot \Big(\frac{n}{w}\Big)^{w/r}.
\]
Since \(w\ge 1\) implies \((n/w)^{w/r}\le n^{w/r}\), we get
\[
\enum_{\Cone,w}\le 2^M\cdot \bigl(e\,n^{1/r}\bigr)^w.
\]
Now \(2^M = 2^{c\log_2 n}=n^c\). For the framework we only use \(w\ge d_{\min}(\Cone)\ge \alpha\log n\), hence
\(n^c \le (2^{c/\alpha})^w\). Folding this polynomial prefactor into the exponential base gives
\[
\enum_{\Cone,w}\le \poly(n)\,\beta^w
\]
for a constant \(\beta>1\) depending only on \(r\) and \(c/\alpha\).
\end{proof}

\paragraph{Remark (stronger forms).}
If one wants the stronger shape \(\enum_{\Cone,w}\le n\,\beta^w\) for all \(w\ge 1\), one can refine the coefficient extraction
and avoid folding \(2^M\) into the base. For our distance proof, the trellis form in Lemma~\ref{lem:outer-dense-spectrum} suffices.

\subsection{Logarithmic minimum distance (O1) via span grouping}
\label{sec:outer-dense-distance}

We now show that a random draw of the taps \(g_t^{(j)}\) yields \(d_{\min}(\Cone)=\Omega(M)=\Omega(\log n)\) with high probability.

\begin{definition}[Span]
Fix a nonzero input \(x\in\F_2^k\). Let \(a\) be the first index with \(x_a=1\) and \(b\) the last index with \(x_b=1\).
Define the span \(\ell:=b-a+1\).
\end{definition}

\begin{lemma}[Weight tail for a fixed span]
\label{lem:outer-span-binomial}
Fix a nonzero input \(x\) of span \(\ell\). Let \(y=\Enc_{\Cone}(x)\) be the outer output under \eqref{eq:outer-dense-encoder}.
Then for each stream \(j\in[r]\), exactly \(\ell+M\) output positions have a nonzero window \((x_t,\dots,x_{t-M})\),
and on those positions \(y_t^{(j)}\) is a uniform bit (independent across \(t\) and across streams).
Consequently,
\[
\wt(y) \;\sim\; \Binomial\bigl(r(\ell+M),\,1/2\bigr),
\]
and for any \(d\ge 1\),
\[
\Pr[\wt(y)\le d]\le (d+1)\binom{r(\ell+M)}{d}\,2^{-r(\ell+M)}
\;\le\;
(d+1)\Big(\frac{e\,r(\ell+M)}{d}\Big)^d 2^{-r(\ell+M)}.
\]
\end{lemma}

\begin{proof}
If the window \((x_t,\dots,x_{t-M})\) is nonzero, then \(\langle g_t^{(j)},(x_t,\dots,x_{t-M})\rangle\) is a nontrivial linear
functional of a uniform vector \(g_t^{(j)}\), hence uniform in \(\F_2\). Independence follows from independence of the tap vectors over
\((t,j)\). The binomial tail bound is a standard union bound over weights \(0,\dots,d\) and \(\binom{L}{d}\le (eL/d)^d\).
\end{proof}

\begin{theorem}[Outer logarithmic minimum distance]
\label{thm:outer-dense-logdistance}
There exists an absolute constant \(\alpha>0\) such that for any choice \(M=\lceil c\log_2 n\rceil\) with \(c\) sufficiently large,
\[
\Pr\bigl[d_{\min}(\Cone)\le \alpha\log_2 n\bigr] = o(1),
\]
where the probability is over the randomness of the taps \(g_t^{(j)}\).
In particular, with high probability \(d_{\min}(\Cone)\ge \alpha\log n\).
\end{theorem}

\begin{proof}
Let \(d:=\alpha\log_2 n\). Union bound over all nonzero inputs \(x\).
Group inputs by span \(\ell\in\{1,\dots,k\}\). For a fixed \(\ell\), there are at most \(k\cdot 2^\ell\le n\cdot 2^\ell\)
nonzero inputs of span \(\ell\) (choose the start position and the \(\ell\)-bit pattern with endpoints \(1\)).
By Lemma~\ref{lem:outer-span-binomial},
\[
\Pr[\wt(\Enc_{\Cone}(x))\le d \mid \text{span}=\ell]
\;\le\;
(d+1)\Big(\frac{e\,r(\ell+M)}{d}\Big)^d 2^{-r(\ell+M)}.
\]
Hence
\[
\Pr[d_{\min}(\Cone)\le d]
\;\le\;
\sum_{\ell=1}^{k} n\,2^\ell\cdot (d+1)\Big(\frac{e\,r(\ell+M)}{d}\Big)^d 2^{-r(\ell+M)}.
\]
Split the sum into two ranges.

\emph{Range A: \(\ell\ge 4d\).}
Here \(\ell+M\ge 4d\), so the factor \(2^\ell\cdot 2^{-r(\ell+M)} \le 2^{-(r-1)\ell}\cdot 2^{-rM}\) decays geometrically in \(\ell\)
(since \(r\ge 2\)). Summing over \(\ell\ge 4d\) contributes at most
\[
n\,(d+1)\Big(\frac{e\,r(\ell+M)}{d}\Big)^d 2^{-rM}\cdot \sum_{\ell\ge 4d} 2^{-(r-1)\ell}
\;\le\;
\poly(\log n)\cdot n^{O(\alpha)}\cdot 2^{-rM}.
\]

\emph{Range B: \(1\le \ell<4d\).}
Here \(\ell+M\le 4d+M\), so we can bound
\[
\Big(\frac{e\,r(\ell+M)}{d}\Big)^d \le \Big(e\,r(4+M/d)\Big)^d = n^{O(\alpha)}
\]
since \(M/d=(c/\alpha)\) is constant.
Also \(\sum_{\ell<4d} 2^\ell \le 2^{4d}=n^{4\alpha}\).
Thus Range B contributes at most
\[
n\cdot 2^{4d}\cdot \poly(\log n)\cdot n^{O(\alpha)}\cdot 2^{-rM}
\;=\;
\poly(\log n)\cdot n^{O(\alpha)}\cdot 2^{-rM}.
\]

Combining, \(\Pr[d_{\min}(\Cone)\le d]\le \poly(\log n)\cdot n^{O(\alpha)}\cdot 2^{-rM}\).
With \(M=c\log_2 n\), we have \(2^{-rM}=n^{-rc}\).
Choosing \(c\) sufficiently larger than the constant exponents induced by \(\alpha\) makes the bound \(o(1)\),
proving the theorem.
\end{proof}

\subsection{Outer interface summary (dense log-memory outer)}
\label{sec:outer-dense-summary}

Combining Theorem~\ref{thm:outer-dense-logdistance} and Lemma~\ref{lem:outer-dense-spectrum}, the dense log-memory outer
satisfies the outer interface used by the serial concatenation framework:
\[
d_{\min}(\Cone)\ge \alpha\log n
\qquad\text{and}\qquad
\enum_{\Cone,w}\le \poly(n)\,\beta^w\ \text{for the relevant weight regime.}
\]
